\vfill\pagebreak
\section{A fixed number of values: arrays}
\label{sec:collections:arrays}

The calendar of Section~\ref{sec:functions:calendar} named the months with a
\token{match} of twelve arms, and counted their days with another. Both are
tables: twelve values, looked up by the number of a month. An \textit{array} is
such a table. It holds a fixed number of values of the same type, its
\textit{elements}, stored one after the other and numbered from 0.

\lstinputlisting[style=coloredverbatim, caption={The months in two arrays, \textit{examples/collections/months.yr}}, label=lst:collections:months]{examples/collections/months.yr}
\vspace{-10pt}%
\begin{lstlisting}[style=bashVerb, escapechar=@+]
@+\textcolor{teal!80}{alice@dev:\mtilde}@+$ ./months
Month: 2
February has 28 days.
\end{lstlisting}

An array is written as its elements between square brackets, separated by
commas, as on lines~4 and~7. A long one can be spread over several lines. Its
type gives the type of the elements and their number, separated by a semicolon:
\token{days} is an \token{[i32 ; 12]}, and \token{names} an
\token{[[c8] ; 12]}, twelve strings. The number is part of the type. An array of
12 elements and one of 13 have different types, and an array never grows or
shrinks. February has 28 days here; the leap years of the calendar are left
out, to keep the listing short.

\subsection{Indices}

\token{days[i]} is the element of \token{days} at \textit{index} \token{i}. The
first element is at index 0, so the last of twelve is at index 11, as
Figure~\ref{fig:collections:array_indices} shows. The user counts the months
from 1, and the array from 0, hence the \token{month - 1} of line~10. The index
can be any integer expression.

\begin{figure}[h]
  \centering
  \begin{adjustbox}{max width=\linewidth}
    \begin{tikzpicture}[cell/.style={rectangle, draw=black!60, fill=background!40,
          minimum width=10mm, minimum height=7mm, font=\ttfamily}]
      \foreach \value [count=\i from 0] in {31, 28, 31, 30, 31, 30, 31, 31, 30, 31, 30, 31} {
        \node[cell] (c\i) at (\i * 1.0, 0) {\value};
        \node[cflab, below=1mm of c\i] {\i};
      }
      \node[cflab, left=3mm of c0] {\token{days}};
      \node[cflab, above=5mm of c1] (feb) {\token{days[1]}};
      \node[cflab, above=5mm of c11] (last) {\token{days[11]}};
      \draw[cfflow] (feb) -- (c1);
      \draw[cfflow] (last) -- (c11);
      \node[cflab, below=6mm of c5.south east] {indices, from 0 to \token{days.len - 1}};
    \end{tikzpicture}
  \end{adjustbox}
  \caption{\label{fig:collections:array_indices} The array \token{days} of
    Listing~\ref{lst:collections:months}. Its twelve elements are stored one
    after the other, and numbered from 0: the days of February are
    \token{days[1]}, those of December \token{days[11]}.}
\end{figure}


The number of elements of an array is its \textit{length}, \token{days.len}. It
is a \token{usize}, the unsigned integer type whose width is that of the
machine (Section~\ref{sec:types:integers}), since a length is never negative.

\subsection{Out of bounds}

\token{days} has no element at index 12. When the index is known while
compiling, the compiler checks it.

\begin{lstlisting}[style=coloredverbatimError, escapechar=@, caption=An index past the end of an array]
use std::io;

fn main() {
  let days = [31, 28, 31, 30, 31, 30, 31, 31, 30, 31, 30, 31];
  println("December has ", @\hb{days[12]}@, " days."); // error
}
\end{lstlisting}

The compiler says \errmsg{E4005}{slice access, index overflow 12us (len) <= 12
  (index)}: the length is 12, and a valid index must be below it.

On line~10 of Listing~\ref{lst:collections:months}, the index depends on what the
user types, and the compiler cannot check it. The program checks it instead,
every time it reads an element, and stops at once when the index is outside the
array.

\begin{lstlisting}[style=bashVerb, escapechar=@+]
@+\textcolor{teal!80}{alice@dev:\mtilde}@+$ ./months
Month: 13
Panic in file "months.yr", at line 10, in function "months::main" !!!
Aborted (core dumped)
\end{lstlisting}

This is a \textit{panic}: the program gives up, and says where. Stopping is the
safe choice. The elements of an array are stored one after the other in memory,
and \token{days[i]} is found by counting \token{i} elements from the first.
Nothing marks the end of the array in memory itself: the element that
\token{days[12]} designates lies right after the last one, where the memory
holds something else, another variable of the program for instance. Without the
check, reading it would give a number that means nothing, and the program would
carry on with it. Writing it would be worse, since it would change that other
variable, which no line of the program names.

A language whose programs cannot touch memory outside what they were given is
said to be \textit{memory safe}. Ymir is, and the check of every index is part
of the price. A program should still check an index that it receives from
outside before using it, to answer with a message rather than a panic;
Section~\ref{sec:collections:ranges} does so for the month.

\notebox{C does not check indices, and writing past the end of an array is the
  \textit{buffer overflow}, one of the most exploited flaws in software. A
  program copies text received from the network into an array of 64
  characters, without checking its length. An attacker sends a longer text,
  whose end overwrites the memory that follows the array, and chooses what is
  written there: among it is the address where the function must go back to
  when it returns. Chapter~\ref{chap:layout} shows what the memory around a
  variable holds, and how the compiler places and aligns the values in it.}

\subsection{Filling an array}

\token{[v ; n]} is an array of \token{n} elements, all equal to \token{v}:
\token{[0 ; 6]} is \token{[0, 0, 0, 0, 0, 0]}. The length must be known while
compiling, since it is part of the type. This form is the usual start of an
array whose elements are computed later, such as the counters of the next
listing, which rolls a die 6\,000 times and counts how often each face comes up.

\lstinputlisting[style=coloredverbatim, caption={Counting the faces of a die, \textit{examples/collections/dice.yr}}, label=lst:collections:dice]{examples/collections/dice.yr}
\vspace{-10pt}%
\begin{lstlisting}[style=bashVerb]
1: 1003
2: 1074
3: 989
4: 945
5: 977
6: 1012
\end{lstlisting}

Each face comes up about 1\,000 times, and the counts change from one run to the
next.

\subsection{Changing the elements}
\label{sec:collections:dmut}

Line~8 of Listing~\ref{lst:collections:dice} adds 1 to one element of
\token{counts}. That is allowed because \token{counts} is declared with
\token{dmut} on line~5. With \token{let} alone, or with \token{let mut}, the
elements of an array cannot change.

\begin{lstlisting}[style=coloredverbatimError, escapechar=@, caption=Changing an element of an array declared with \token{let}]
use std::io;
use std::rand;

fn main() {
  let counts = [0 ; 6];
  for _ in 0 .. 6000 {
    let die = uniform(1, 6);
    @\hb{counts[die - 1]}@ += 1; // error
  }
  println(counts);
}
\end{lstlisting}

The compiler says \errmsg{E4082}{left operand of type i32 is immutable}: the
element, an \token{i32}, is a constant. A collection has two things that can
change: the variable, which can be given a whole new collection, and the
elements inside it. The declaration says which ones may.

\begin{center}
  \begin{adjustbox}{max width=\linewidth}
    \begin{tabular}{|l|l|}
      \hline
      Declaration & What may change\\[0pt]
      \hline
      \hline
      \token{let counts = [0 ; 6];} & nothing\\[0pt]
      \token{let mut counts = [0 ; 6];} & the variable: \token{counts = [1, 1, 1, 1, 1, 1];} is allowed\\[0pt]
      \token{let dmut counts = [0 ; 6];} & the variable and its elements: \token{counts[0] = 5;} is allowed too\\[0pt]
      \hline
    \end{tabular}
  \end{adjustbox}
\end{center}

The \token{d} of \token{dmut} stands for \textit{deep}: the mutability goes all
the way down, to the elements. The same rule applies to the fields of a tuple,
so \token{pair.0 = 5;} needs a \token{pair} declared with \token{dmut}.
Why Ymir tells the two levels apart only shows when several variables refer to
the same elements, which the next chapter is about. Until then, the rule is
enough: a collection whose elements change is declared \token{dmut}.

\subsection{Going through an array}

\token{for x in a} runs its block once for each element of \token{a}, in order
of index, and \token{x} holds the current element. With two names, as on
line~10 of Listing~\ref{lst:collections:dice}, the first holds the index, a
\token{usize}, and the second the element. Line~11 prints \token{i + 1}, so that
the faces are numbered from 1, as on the die.

Both names are constants, like the iterator of a range: \token{count = 0;} in
the block would not reset an element of \token{counts}, and is refused. To
change the elements, the block goes through the indices and writes
\token{counts[i]}, or uses the loops of the next chapter.

\subsection{Expanding an array}
\label{sec:collections:expand_array}

\token{expand} (Section~\ref{sec:collections:expand}) replaces a tuple with its
fields, which requires the compiler to know how many there are. The length of
an array is part of its type, so the compiler knows it too, and an array can be
expanded like a tuple.

\lstinputlisting[style=coloredverbatim, caption={Expanding an array, \textit{examples/collections/seasons.yr}}, label=lst:collections:seasons]{examples/collections/seasons.yr}
\vspace{-10pt}%
\begin{lstlisting}[style=bashVerb]
24
[3, 4, 5, 6, 7, 8], 6 months
(3, 4, 5, spring)
\end{lstlisting}

\begin{itemize}
\item \textit{In a call.} On line~9, \token{volume(expand size)} is the call
  \token{volume(2, 3, 4)}: each element of the array is an argument.
\item \textit{In an array.} On line~13, \token{[expand spring, expand summer]}
  is the array \token{[3, 4, 5, 6, 7, 8]}, the elements of \token{spring}
  followed by those of \token{summer}. Its type is \token{[i32 ; 6]}: the
  compiler adds the lengths.
\item \textit{In a tuple.} On line~16, \token{(expand spring, "spring")} is a
  tuple of four fields, three numbers and a string.
\end{itemize}

The slices of the next section have a length known only when the program runs,
in general, and cannot be expanded whole: a slice received as a parameter is
refused with \errmsg{E4257}{unknown length of expansion for type [i32]}. A part
of a slice whose bounds are known while compiling can be
(Section~\ref{sec:collections:expand_slice}).

\subsection{Arrays of arrays}

The elements of an array can be of any type, arrays included.
\token{[[0 ; 3] ; 3]} is a grid of 3 rows of 3 zeros, of type
\token{[[i32 ; 3] ; 3]}. \token{grid[1]} is its second row, and
\token{grid[1][2]} the third element of that row.

\begin{lstlisting}[style=coloredverbatim]
let dmut grid = [[0 ; 3] ; 3];
grid[1][2] = 5;
println(grid[1]);
for row in grid {
  for cell in row {
    print(cell, " ");
  }
  println();
}
\end{lstlisting}
\vspace{-10pt}%
\begin{lstlisting}[style=bashVerb]
[0, 0, 5]
0 0 0
0 0 5
0 0 0
\end{lstlisting}

\token{grid[1][2] = 5;} changes one element of one row, which needs
\token{grid} declared \token{dmut}, as for any element. The loops go through
the grid in two levels: the outer loop gives each row, an array of 3 elements,
and the inner loop each element of that row.

Exercise~9 plays on such a grid, and Section~\ref{sec:collections:grids} comes
back to grids whose size is known only at run time.
